\chapter{Abuse}

\section*{Definition and Overview}
Abuse is the deliberate mistreatment, neglect, or exploitation of another person, usually by someone in a position of power, trust, or care. It takes many forms, including physical, emotional (psychological), sexual, and financial abuse, together with neglect, which is the failure to provide the care and support a person needs. Abuse can affect anyone at any stage of life, but it most often involves children, intimate partners, and older people. It is a serious health problem because of the direct injuries it causes and the profound long-term physical and mental health consequences that frequently follow.

\section*{Epidemiology}
Abuse is common and substantially under-reported. Approximately one in four women and one in ten men experience intimate partner violence at some point in their lives, and an estimated one in six older people experiences some form of abuse in a community setting each year. Child abuse and neglect affect a significant number of children annually, and many victims never disclose what has happened. Abuse crosses every age, income, culture, and level of education, although poverty, social isolation, disability, family stress, and substance misuse increase the risk.

\section*{Aetiology and Pathophysiology}
Abuse arises from a combination of individual, relationship, family, and social factors. A common thread is an imbalance of power in which one person controls or dominates another. Contributing factors include substance misuse, mental illness, high stress, poverty, social isolation, a personal history of having been abused, and cultural attitudes that accept violence. Children are made more vulnerable by young age and disability, while in older people, dependence on a carer is a key risk factor.

\begin{itemize}
    \item \textbf{Physical abuse:} direct injury from hitting, shaking, burning, or the use of restraint
    \item \textbf{Emotional abuse:} humiliation, threats, isolation, and controlling behaviour
    \item \textbf{Sexual abuse:} any sexual act without consent, including against children
    \item \textbf{Financial abuse:} stealing money or property, or controlling a person's finances without consent
    \item \textbf{Neglect:} failure to provide food, shelter, medical care, or adequate supervision
\end{itemize}

The mechanism of harm is twofold. There is direct physical injury, and there is the effect of chronic stress: prolonged exposure to threat and fear alters stress hormone systems and is thought to underlie the higher rates of depression, anxiety, cardiovascular disease, and other chronic conditions seen in people who have experienced abuse.

\section*{Clinical Features}
The signs of abuse vary with the type and the victim, but certain features should raise concern.

\begin{itemize}
    \item Injuries that do not match the explanation given, or that appear in clusters, on soft tissue areas, or in the shape of an object
    \item Bruises at different stages of healing, bite marks, or burns with a clear outline
    \item Delays between injury and seeking care, or frequent visits to different clinicians
    \item Fear, withdrawal, anxiety, or flinching when approached
    \item Changes in behaviour, sleep, or school performance in children
    \item Poor hygiene, inappropriate clothing, hunger, or signs that medical needs are being neglected
    \item In older people: fear of a particular carer, unexplained bruising, pressure injuries, or sudden financial changes
    \item Depression, post-traumatic stress, substance misuse, and self-harm
\end{itemize}

\section*{Differential Diagnosis}
Because abuse can be mistaken for other conditions, and vice versa, careful assessment is essential.

\begin{itemize}
    \item Accidental injury, including from play, sport, or falls
    \item Bleeding disorders that cause easy bruising, such as low platelet levels or haemophilia
    \item Medical conditions causing unexplained bruising or bone injuries in children, including vitamin deficiencies and bone disease
    \item Cultural practices that produce marks or bruising, such as coining or cupping
    \item Self-harm or suicide attempts
    \item Skin conditions that leave marks resembling injury
\end{itemize}

\section*{When to Seek Medical Care}
Anyone who is being abused, or who suspects abuse of someone they know, should tell a doctor, nurse, social worker, teacher, or police officer. Health workers can assist with safety planning, treatment of injuries, and referral to specialist services.

\textbf{Call 000 (or local emergency services) for:}
\begin{itemize}
    \item Immediate danger of being hurt, or a serious injury
    \item A child who is at risk of harm
    \item Severe bleeding, head injury, or signs of recent sexual assault
    \item A person who is in danger from an abuser at that moment
\end{itemize}

\section*{Diagnosis and Investigations}
Recognition and documentation require care, sensitivity, and a structured approach.

\begin{itemize}
    \item \textbf{History:} where possible the person is seen alone and asked directly, in a private and supportive setting
    \item \textbf{Full examination:} careful documentation of injuries, including bruises, fractures, and other marks, using body diagrams and photography where appropriate
    \item \textbf{Imaging:} a skeletal survey in young children to look for old or healing fractures
    \item \textbf{Blood tests:} to exclude bleeding disorders when bruising is unexplained
    \item \textbf{Forensic examination:} by a trained team after sexual assault, including testing for sexually transmitted infections and, where indicated, pregnancy
    \item \textbf{Psychosocial assessment:} of immediate safety, mental health, and the person's support network
    \item \textbf{Reporting:} health workers are legally required in most jurisdictions to report suspected abuse of children and vulnerable adults
\end{itemize}

\section*{Management}
Management places safety first and is delivered in a respectful, non-judgemental manner.

\begin{itemize}
    \item \textbf{Safety first:} safety planning, urgent protection, and assistance to reach a safe place such as a refuge when needed
    \item \textbf{Medical treatment:} care of injuries, and emergency contraception and infection prevention after sexual assault
    \item \textbf{Reporting:} involvement of child protection authorities or police where the law requires it or where children are at risk
    \item \textbf{Trauma-informed care:} an approach that recognises the effects of trauma and avoids re-traumatising the person
    \item \textbf{Psychological support:} counselling and therapy for post-traumatic stress, depression, and anxiety
    \item \textbf{Practical support:} housing, legal protection such as intervention orders, financial assistance, and advocacy
    \item \textbf{Follow-up:} ongoing medical, mental health, and social support while the person rebuilds their life
\end{itemize}

\section*{Complications}
\begin{itemize}
    \item Physical injury, disability, and, in the worst cases, death
    \item Depression, anxiety, post-traumatic stress disorder, and suicide
    \item Substance misuse as a coping strategy
    \item Chronic pain, cardiovascular disease, and other stress-related illness
    \item In children: developmental delay, behavioural problems, and difficulty forming trusting relationships
    \item Intergenerational harm, in which people abused as children are more likely to experience or perpetrate violence later in life
    \item Homelessness and poverty, especially for people leaving an abusive relationship
\end{itemize}

\section*{Prognosis and Outcomes}
With safety, support, and time, many people who have experienced abuse recover well. Removing a child from harm and providing stable, caring environments leads to markedly better outcomes, and psychological therapy can significantly reduce post-traumatic symptoms. Without intervention, abuse often continues and may escalate, and the physical and mental health consequences tend to accumulate. Recovery is a gradual process and varies greatly from person to person, but early recognition and coordinated support improve the outlook considerably.

\section*{Prevention and Self-Care}
\begin{itemize}
    \item Learn about healthy, respectful relationships and teach children about body safety and consent
    \item Support early childhood programs and parenting help that reduce the stresses linked to abuse
    \item Screen for abuse routinely in health care so that people can be asked privately and safely
    \item If you are in an abusive relationship, plan for safety: keep important documents and money accessible and have a plan for leaving
    \item Reach out to trusted people and to specialist services rather than coping alone
    \item Address alcohol and drug problems and seek help for mental illness within the family
    \item Be alert to signs of abuse in older or dependent relatives, and speak up if you are concerned
\end{itemize}

\section*{Sources and Further Reading}
The following organisations provide reliable information and support resources on abuse and its effects.

\begin{itemize}
    \item World Health Organization. \textit{Violence prevention and child maltreatment resources.} https://www.who.int/
    \item Centers for Disease Control and Prevention. \textit{Intimate partner violence prevention.} https://www.cdc.gov/
    \item NHS. \textit{Abuse and domestic violence: where to get help.} https://www.nhs.uk/live-well/abuse-and-domestic-violence/
    \item Mayo Clinic. \textit{Domestic violence against women: overview and support.} https://www.mayoclinic.org/diseases-conditions/domestic-violence/
    \item MedlinePlus. \textit{Child abuse and domestic violence information.} https://medlineplus.gov/childabuse.html
    \item Australian Government Department of Social Services. \textit{Family and domestic violence support services.} https://www.dss.gov.au/
\end{itemize}
